\documentclass[11pt]{article}
\usepackage[a4paper,margin=1in]{geometry}
\usepackage{amsmath,amssymb,amsthm,booktabs,microtype}
\newtheorem{proposition}{Proposition}[section]
\theoremstyle{remark}\newtheorem{remark}[proposition]{Remark}
\title{A Typed Ledger for Noisy Spectral Tails, Graded Counterloops, and Their Missing Glue}
\author{Bin Wang}\date{July 2026}
\begin{document}\maketitle
\begin{abstract}
The variable-rank spectral tail of RH-282 fills one entry of the noisy branch,
but it cannot be transferred silently to the monodromy counterloop branch.
We update the two five-obligation ledgers and add an explicit cross-branch
glue bit.  The noisy modulus-spectral vector is $(1,0,1,1,1)$; the graded
monodromy vector is $(1,1,0,1,1)$; the weighted glue bit is zero.  Each branch
therefore has score four, with complete count zero.  Their coordinatewise
maximum is the all-ones vector, but that merge is ill typed: the spectral
head and graded head are not identified.  RH-288 states the exact weighted
complement-to-anchor interface needed to legalize the merge.  A sufficient
decomposition requires both a weighted total-trace bridge and weighted
head-to-counterloop transport.  Gates A--E remain open.
\end{abstract}

\section{Typed obligations}
The five entries are:
\begin{enumerate}
 \item a legal head;
 \item a coefficient bridge to the deterministic numerator;
 \item a uniform high-order tail;
 \item an analytic target tail;
 \item a certified target boundary constant.
\end{enumerate}
The same labels are used for both branches, but the underlying heads are
different objects.  The noisy branch uses the actual modulus-complete
algebraic eigenvalue multiset.  The graded branch uses the deterministic
finite-radius monodromy counterloop.

\begin{center}\small
\begin{tabular}{lccccc|c}\toprule
branch & head & bridge & tail & target & boundary & score\\\midrule
noisy modulus spectrum & 1&0&1&1&1&4\\
graded monodromy counterloop & 1&1&0&1&1&4\\\bottomrule
\end{tabular}
\end{center}

\section{No coordinatewise merge}
\begin{proposition}[typed-union obstruction]
The coordinatewise maximum of the two vectors is $(1,1,1,1,1)$, but it is
not a complete certificate unless the noisy modulus complement satisfies the
weighted prefix bridge to the deterministic anchor.  By RH-288, it is
sufficient to prove on the same clock both the weighted total-noisy-trace
bridge to counterloop plus anchor and the weighted noisy-head bridge to the
counterloop.
\end{proposition}
\begin{proof}
The spectral tail estimate controls the complement of the modulus-selected
noisy eigenvalues.  The graded coefficient bridge subtracts moments of a
different finite multiset.  Without an equality or asymptotically vanishing
weighted complement-to-anchor prefix, the tail of one factor and the prefix
of the other do not occur in a common determinant decomposition.  Writing
total trace as head plus complement, RH-288 decomposes that missing prefix
into a total-trace/counterloop/anchor error minus a head/counterloop error.
Controlling only the latter leaves the former open.  The compound interface
is currently false/open, so coordinatewise union changes the object
mid-proof.
\end{proof}

\section{Status}
The cross-branch weighted-glue bit is zero.  Both branch-complete counts are
therefore zero even though both raw scores equal four.  RH-286 improves the
finite-radius diagnostic target and RH-287 gives a rate-free growing prefix;
neither supplies the two synchronized weighted estimates or the direct
complement-to-anchor prefix.

\begin{remark}
No Gate-A determinant identification has been completed.  Gates B--E also
remain false/open.  Nothing here constructs a Hilbert--Polya operator,
identifies Riemann zeros, proves a von Mangoldt trace identity, or proves RH.
\end{remark}
\end{document}
